\section{Random Features: A Case Study}
\label{sec:rff}
\begin{figure}[H]
    \centering
    \includegraphics[width=\textwidth]{rff_fmnist.png} \\
    %\includegraphics[width=0.3\textwidth]{tmp_figures/rff_fmnist_slice_sam.png}
    %\includegraphics[width=0.3\textwidth]{tmp_figures/rff_fmnist_slice_emb.png}
    \caption{\textbf{Random Fourier Features} on the Fashion MNIST dataset. The setting is equivalent to two-layer neural network with $e^{-ix}$ activation, with randomly-initialized first layer that is fixed throughout training. The second layer is trained using gradient flow.}
    \label{fig:rff}
\end{figure}


In this section, for completeness sake, we show that both the model- and sample-wise double descent phenomena are not unique to deep neural networks---they exist even in the setting of Random Fourier Features of \cite{rahimi2008random}. This setting is equivalent to a two-layer neural network with $e^{-ix}$ activation. The first layer is initialized with a $\mathcal{N}(0, \frac{1}{d})$ Gaussian distribution  and then fixed throughout training. The width (or embedding dimension) $d$ of the first layer parameterizes the model size. The second layer is initialized with $0$s and trained with MSE loss.

Figure \ref{fig:rff} shows the grid of Test Error as a function of
both number of samples $n$ and model size $d$.
Note that in this setting $\EMC = d$ (the embedding dimension). As a result, as demonstrated in the figure, the peak follows the path of $n=d$. Both model-wise and sample-wise (see figure~\ref{fig:rff-samples})  double descent phenomena are captured, by horizontally and vertically crossing the grid, respectively. 

\begin{figure}[h]
    \centering
    \includegraphics[width=0.6\textwidth]{figures/rff_fmnist_samples.png} \\
    %\includegraphics[width=0.3\textwidth]{tmp_figures/rff_fmnist_slice_sam.png}
    %\includegraphics[width=0.3\textwidth]{tmp_figures/rff_fmnist_slice_emb.png}
    \caption{Sample-wise double-descent slice for Random Fourier Features on the Fashion MNIST dataset. In this figure the embedding dimension (number of random features) is 1000.}
    \label{fig:rff-samples}
\end{figure}